\documentclass{article}
\usepackage{enumerate}
\usepackage[margin=1.4in]{geometry}
\title{Teaching Statement}
\author{Joel Villatoro}
\date{}


\begin{document}
\vspace{-0.5in}
\maketitle


I consider teaching and mentorship one of my primary responsibilities as an academic mathematician. I also consider it very rewarding part of my job.




My teaching philosophy has been shaped by my experiences so in order to explain my thoughts on teaching it will be important for me to describe my various experiences and the conclusions I have drawn from them. Before doing so, however, I can at least summarize my conclusions with the following opinion: Students generally learn very little by passive absorption and memorization without reflection. However, it is very easy for both teachers and a students to fall into a pattern of interaction that has very little engagement. As a teacher, it is critical to use whatever tools are at your disposal to maximize the amount of time, within reason, that your students spend actively thinking about the content of your course.

I have been fortunate to teach courses with very different student body compositions as well as methodologies. 
At the University of Illinois, teaching assistants for large lecture courses were encouraged to provide their students with an active learning style of discussion sections. At KU Leuven, teaching assistants play a similar role but students are largely expected to complete exercises on their own. 
This change in methodology was a challenge for me. 
By this point I had realized the importance of maximizing student engagement but the tools at my disposal for achieving that changed quite suddenly. 
I realized that at KU Leuven my role as a teaching assistant became much more important, in a sense. 
I was the only person in the room with whom they were expected to interact with. 
Even more challenging, I observed that the culture of the students at KU Leuven was much more reserved. Belgian students seemed to be less likely to ask for assistance and were often reluctant to interact with anyone at all. To mitigate these obstructions to learning, I was forced to adapt in a few ways. One adaptation was to prod students with questions such as ``How is the worksheet going?" much more frequently. This was often successful but some students were still often reluctant to engage or were embarrassed to be seen by their peers as asking `easy questions.' In such cases, I realized that I could still indirectly interact with such students by looking over their shoulder and doing example problems in front of the entire class which were similar to the ones they were struggling with.

During the Spring of 2020 I was tasked with teaching a graduate course in Symplectic Geometry. My original plan was to closely emulate the historical format for previous instances of the same course. This meant that it was originally going to be standard lecture course with two graded assignments as well an oral exam. The main deviation I undertook was to include as many exercises as possible in the lecture notes. My intuition was that graduate students would be generally more motivated and that the best way to maximize their engagement was to provide them with as many opportunities as possible. I was correct that graduate students were more motivated, but in the future I think that the students would have received an additional boost by giving them more graded assignments with feedback.

Halfway through the course in Symplectic Geometry, the rising COVID pandemic meant that changes were going to be necessary to ensure the health and safety of both students and instructors. The class was switched to an asynchronous lecture format punctuated with live question and answer sessions. It was a challenge to adapt my teaching style to the new format and I spent a lot of time learning the ins and outs of producing prerecorded video lectures as well as managing online question and answer sessions. My primary goal in distance learning environments is to do the best I can to encourage the students to consume the class material on a regular schedule. One way to do this is to combine the video lectures with live problem solving sessions, preferably with mandatory (graded) attendance.

Another important component of my teaching is mentorship. One major advantage to teaching in a mentor capacity is that there are comparatively few barriers to student engagement. For this reason, I have often felt that the mentorship roles I played were the most satisfying. For this reason, I like to seek out find mentorship roles to play in parallel with my more traditional teaching. At the University of Illinois I did this by mentoring undergraduate students via the Illinois Geometry Lab as well as graduate students through the SLOAN peer mentorship program. The second example is of particular importance to me since it provided me with the opportunity to be a Latino-American role-model for a younger graduate student with a similar background. As a postdoc at KU Leuven, I realized that an important part of my role in the research group was to act as a collaborator and mentor with the graduate students. In addition to being satisfying, this work was ultimately very fruitful as it resulted in a paper which I coauthored with a graduate student.

In light of the increasing prevalence of distance learning classroom formats, finding new strategies to ensure that students remain engaged with their class material is more important than ever. I look forward to continuing to grow as a teacher as I learn new strategies to share my knowledge and expertise with the next generation of mathematicians.

\end{document}
